%% =============================================================================
%% DISCUSSION
%% =============================================================================

\section{Discussion} \label{sec:discussion}

This paper described chat-first data analytics---separating
interpretation from computation, pre-validating methods rather than
generating code, and keeping data local---and implemented the approach
in \pkg{prompt2analytics}, a \proglang{Rust} library of over 200
validated methods exposed through MCP. The evaluation
(Section~\ref{sec:evaluation}) shows that tool selection works well:
90--100\% accuracy on curated prompts, with errors concentrated in
edge cases requiring fine distinctions between related methods.
Several limitations qualify these results. The practical benefits of
chat-first interfaces---accessibility, reduced cognitive load,
exploratory dialogue---remain hypothesized, as no user studies were
conducted. Controlled experiments comparing task completion time, error
rates, and user satisfaction against traditional code-based workflows
are needed and represent the most important direction for future
research. When the correct tool is selected, manual inspection of
recorded tool calls shows that core parameters are almost always
correct; the main imprecision is omission of optional arguments that
fall back on sensible defaults. Users should verify that the system
ran the analysis they intended, not merely that it selected the right
method. While method coverage is broad, important categories
remain unimplemented---Bayesian inference, Heckman selection, Conley
spatial standard errors, and structural equation modeling among
them (Table~\ref{tab:method-coverage} in
Appendix~\ref{app:best-practices} documents gaps and recommends
alternatives). Data are loaded entirely into memory, limiting
applicability to datasets that fit in RAM. And because the MCP schema
uses semantic versioning (currently 0.x), the API may change as the
software matures.  We follow a conservative deprecation policy: tool
names and required parameters will not be removed or renamed without
a major version increment; new optional parameters can be added in
minor releases.  Tool call transcripts from the examples in this
paper will remain valid in all 0.x releases.  The underlying MCP
protocol (\pkg{rmcp}~0.8) is itself pre-1.0; protocol-level changes
are absorbed by the server without affecting tool schemas.

These limitations suggest clear directions for future work. Reducing
omission of optional parameters---through structured output
constraints, two-pass verification where the model checks its own
extraction before execution, or richer schema descriptions that make
optional parameters more salient---would narrow the gap between
correct tool selection and fully specified analysis execution.
Property-based testing could strengthen correctness guarantees beyond
comparison to reference implementations; QuickCheck-style testing for
statistical properties (e.g., residuals orthogonal to regressors) would
complement numerical validation. Foreign function interfaces to
\proglang{R} and \proglang{Python} would let users call the
\proglang{Rust} implementations from familiar environments, combining
compiled performance with existing analytical workflows. An
alternative design would have been to build \pkg{p2a-core} as an
\proglang{R} package from the start, leveraging \proglang{R}'s
ecosystem and user base. We chose a standalone architecture because the
primary interface is the MCP server---a long-running process that any
LLM client can connect to---which does not fit naturally into
\proglang{R}'s package model. FFI wrappers remain a priority for
giving \proglang{R} and \proglang{Python} users direct access to the
compiled implementations. If extraction
quality improves sufficiently, the bottleneck shifts from ``did the
model understand the request?'' to ``does the method library cover this
analysis?''---at which point extending coverage becomes more valuable
than improving the language model, which is exactly the division of
labor the architecture is designed to exploit.

Even with these extensions, the system's role would remain
complementary. \pkg{prompt2analytics} is not a replacement for \proglang{R},
\proglang{Stata}, or \proglang{Python}. It occupies a different niche:
users who know what analysis they want but not which software command
produces it. A researcher fluent in \proglang{R} has no need for a
conversational interface to OLS; a policy analyst who has never used
\proglang{R} might. The system complements existing tools rather than
competing with them---and for users who outgrow the conversational
interface, every tool call maps to a public function that can be
called directly from \proglang{Rust} or, once foreign function
interfaces are available, from \proglang{R} and \proglang{Python}.

Maintaining a library of over 200 validated methods is a legitimate
concern. The same AI-assisted development tools that motivate the
project also address this: coding agents (e.g., Claude
Code) can implement new methods, write validation
tests, update dependencies, and respond to bug reports filed as
GitHub issues---with human review of each pull request ensuring
correctness. This workflow already produced the majority of the
current implementation. It scales sublinearly with method count
because new methods reuse established patterns (the
\code{LinearEstimator} trait, shared matrix operations, the MCP
handler template), and because validation against \proglang{R}
reference implementations is itself automatable. The project is
open-source under the MIT license; community contributions follow
the same review process.

More broadly, the chat-first approach exemplifies a pattern applicable
beyond econometrics: using language models as interpreters between human
intent and specialized software~\citep{Yao+etal:2023,Schick+etal:2023}.
Any domain where users express goals in natural language and systems
execute well-defined operations---laboratory automation, geographic
information systems, circuit simulation---could adopt the same
architecture, particularly as structured protocols like
MCP~\citep{MCP:2024} standardize the interface between models and tools.

The LLM landscape evolves rapidly.  Models that achieve 100\% on the
current 87-test benchmark (Claude~3.5 Haiku, Qwen~2.5 72B) suggest the
benchmark may already exhibit ceiling effects; expanding it with
harder multi-turn and naturalistic cases is a priority.  The
evaluation harness and test cases are open-source, so new models can
be benchmarked as they appear.  Importantly, improving LLM capability
does not undermine the architecture---it strengthens it.  Better
models reduce parameter omissions and may eventually make
code-generation approaches more reliable, but the
benefits of pre-validation and data locality are orthogonal to model
quality.  Even if future models generate flawless code, users who
require auditable, reproducible analyses on local data would still
benefit from fixed, validated tool implementations.

The success of this approach depends on a clear division of labor:
language models interpret and explain; specialized software computes.
Attempts to have language models do both---generating code that performs
numerical computation---conflate interpretation with execution and
sacrifice the reliability that scientific work
requires~\citep{Schick+etal:2023}. The viability of local LLM deployment shown in
Section~\ref{sec:deployment} has implications beyond data privacy.
If small models can reliably select among pre-implemented methods,
conversational interfaces to scientific software become possible
without cloud connectivity or high-end hardware. The design of
\pkg{prompt2analytics} suggests a template: implement algorithms in
high-performance, memory-safe languages; expose them through structured
protocols; and let language models handle the translation between human
intent and computational execution.


%% =============================================================================
%% ACKNOWLEDGMENTS
%% =============================================================================

\section*{Acknowledgments}

The authors thank the reviewers for constructive feedback that improved
the paper.


%% =============================================================================
%% APPENDIX
%% =============================================================================

\clearpage
\appendix

%% Renumber tables and figures by appendix section
\renewcommand{\thetable}{\Alph{section}.\arabic{table}}
\renewcommand{\thefigure}{\Alph{section}.\arabic{figure}}

\section*{Appendix}
\addcontentsline{toc}{section}{Appendix}

\vspace{1em}
\noindent\textbf{Contents}
\vspace{0.5em}

\noindent\hyperref[app:best-practices]{A\hspace{1em}Best practices and coverage gaps}\\
\noindent\hyperref[app:llm-eval]{B\hspace{1em}LLM evaluation details}\\
\noindent\hyperref[app:latency]{C\hspace{1em}API response latency}\\
\noindent\hyperref[app:deployment]{D\hspace{1em}Deployment guide}\\
\noindent\hyperref[app:privacy]{E\hspace{1em}Privacy and security}\\
\noindent\hyperref[app:reproducibility]{F\hspace{1em}Reproducibility}\\
\noindent\hyperref[app:benchmark-methodology]{G\hspace{1em}Benchmark methodology}\\
\noindent\hyperref[app:implementation]{H\hspace{1em}Implementation notes}\\
\noindent\hyperref[app:extended-examples]{I\hspace{1em}Extended examples}

\vspace{2em}
